\documentclass[aspectratio=169, 14pt]{beamer}
\usepackage{stampfly_slides}

\lesson{6}{PID 制御}{PID Control}

\begin{document}

% ------------------------------------------------------------------
% Slide 1: Title
% ------------------------------------------------------------------
\begin{frame}
  \titlepage
\end{frame}

% ------------------------------------------------------------------
% Slide 2: Goal (updated)
% ------------------------------------------------------------------
\begin{frame}{今日のゴール / Today's Goal}
  \begin{block}{目標}
    工学形式 PID を実装し、不完全微分でノイズに強い制御を実現する
  \end{block}

  \vspace{0.8em}

  \begin{itemize}
    \item P 制御の限界（定常偏差、振動）
    \item 工学形式 $K_p / T_i / T_d$ とは
    \item 不完全微分: ノイズに強い微分項
    \item ログから調整する方法
  \end{itemize}
\end{frame}

% ------------------------------------------------------------------
% Slide 3: PID Block Diagram (TikZ updated)
% ------------------------------------------------------------------
\begin{frame}{PID 制御器 / PID Controller}
  \begin{center}
    \resizebox{0.78\textwidth}{!}{\documentclass[border=10pt]{standalone}
\usepackage{tikz}
\usepackage{amsmath}
\usetikzlibrary{arrows.meta, positioning, shapes.geometric, calc}

\begin{document}

\definecolor{pcolor}{RGB}{0, 120, 200}    % Blue - P
\definecolor{icolor}{RGB}{40, 160, 40}    % Green - I
\definecolor{dcolor}{RGB}{220, 50, 30}    % Red - D
\definecolor{arrowcolor}{RGB}{80, 80, 80}
\definecolor{sumcolor}{RGB}{120, 80, 150} % Purple - sum
\begin{tikzpicture}[
    block/.style={
      rectangle, draw=#1, line width=1.5pt,
      fill=#1!8, rounded corners=3pt,
      minimum width=35mm, minimum height=13mm,
      font=\bfseries, text=#1, align=center
    },
    sum/.style={
      circle, draw=sumcolor, line width=1.5pt,
      fill=sumcolor!8, minimum size=10mm,
      font=\large\bfseries, text=sumcolor
    },
    arrow/.style={-{Stealth[length=3mm]}, line width=1.2pt, arrowcolor},
  ]

  % Input
  \node[font=\bfseries, arrowcolor] (input) at (-1, 0) {error $e$};

  % Branch point
  \coordinate (branch) at (1.5, 0);

  % P block (top)
  \node[block=pcolor] (P) at (5.5, 2) {P: $K_p \cdot e$};

  % I block (center)
  \node[block=icolor] (I) at (5.5, 0) {I: $K_i \displaystyle\int e\,dt$};

  % D block (bottom)
  \node[block=dcolor] (D) at (5.5, -2) {D: $K_d \dfrac{de}{dt}$};

  % Sum
  \node[sum] (sum) at (10, 0) {$\Sigma$};

  % Output
  \node[font=\bfseries, arrowcolor] (output) at (12.5, 0) {output $u$};

  % Input arrow to branch
  \draw[arrow] (input) -- (branch);

  % Branch to P/I/D
  \draw[arrow, pcolor] (branch) |- (P);
  \draw[arrow, icolor] (branch) -- (I);
  \draw[arrow, dcolor] (branch) |- (D);

  % P/I/D to sum
  \draw[arrow, pcolor] (P) -| (sum);
  \draw[arrow, icolor] (I) -- (sum);
  \draw[arrow, dcolor] (D) -| (sum);

  % Sum to output
  \draw[arrow] (sum) -- (output);

  % Plus signs at sum
  \node[font=\scriptsize, pcolor] at ($(sum)+(120:0.75)$) {$+$};
  \node[font=\scriptsize, icolor] at ($(sum)+(180:0.75)$) {$+$};
  \node[font=\scriptsize, dcolor] at ($(sum)+(240:0.75)$) {$+$};

  % Labels (left-aligned just right of blocks, above/below arrow lines)
  \node[font=\small, pcolor, anchor=south west] at (7.5, 2) {Proportional};
  \node[font=\small, icolor, anchor=south west] at (7.5, 0) {Integral};
  \node[font=\small, dcolor, anchor=north west] at (7.5, -2) {Derivative};

\end{tikzpicture}
\end{document}
}
  \end{center}
\end{frame}

% ------------------------------------------------------------------
% Slide 4: Why Incomplete Derivative? (new)
% ------------------------------------------------------------------
\begin{frame}{なぜ不完全微分か / Why Incomplete Derivative?}
  \begin{columns}[T]
    \begin{column}{0.48\textwidth}
      \begin{alertblock}{問題}
        \begin{itemize}
          \item ジャイロにはノイズがある
          \item 理想微分 $\frac{de}{dt}$ はノイズを増幅
          \item 高周波振動 $\to$ モータに悪影響
        \end{itemize}
      \end{alertblock}
    \end{column}
    \begin{column}{0.48\textwidth}
      \begin{exampleblock}{解決策}
        \begin{itemize}
          \item 不完全微分フィルタで高周波カット
          \item $\eta = 0.1 \sim 0.2$
          \item 小 $\to$ フィルタ弱い
          \item 大 $\to$ フィルタ強い
        \end{itemize}
      \end{exampleblock}
    \end{column}
  \end{columns}

  \vspace{0.8em}

  \begin{block}{トレードオフ}
    微分の効き（応答速度）$\longleftrightarrow$ ノイズ除去（安定性）
  \end{block}
\end{frame}

% ------------------------------------------------------------------
% Slide 5: Engineering Form Parameters (new)
% ------------------------------------------------------------------
\begin{frame}{工学形式パラメータ / Engineering Form}
  \begin{table}
    \centering
    \footnotesize
    \begin{tabular}{llp{55mm}}
      \hline
      \textbf{パラメータ} & \textbf{意味} & \textbf{効果} \\
      \hline
      $K_p$  & 比例ゲイン   & 大$\to$応答速い、大きすぎ$\to$振動 \\
      $T_i$  & 積分時間 [s] & 小$\to$積分強い（偏差早く消える）、小さすぎ$\to$ワインドアップ \\
      $T_d$  & 微分時間 [s] & 大$\to$微分強い（振動抑制）、大きすぎ$\to$ノイズ増幅 \\
      $\eta$ & フィルタ係数  & $0.1 \sim 0.2$、大$\to$ノイズに強い \\
      \hline
    \end{tabular}
  \end{table}

  \vspace{0.5em}

  \begin{block}{教科書形式との対応}
    $K_i = K_p / T_i$\,,\quad $K_d = K_p \cdot T_d$
  \end{block}
\end{frame}

% ------------------------------------------------------------------
% Slide 6: Log-Based Tuning Guide (new -- participant request)
% ------------------------------------------------------------------
\begin{frame}{ログから調整する / Log-Based Tuning}
  \begin{table}
    \centering
    \footnotesize
    \begin{tabular}{lll}
      \hline
      \textbf{ログで見える症状} & \textbf{原因} & \textbf{調整} \\
      \hline
      振動が収まらない   & $K_p$ 過大          & $K_p\!\downarrow$ \\
      応答が遅い         & $K_p$ 過小          & $K_p\!\uparrow$ \\
      定常偏差が残る     & I 項不足            & $T_i\!\downarrow$（積分強化）\\
      オーバーシュート大 & D 項不足            & $T_d\!\uparrow$（微分強化）\\
      高周波ノイズ       & $\eta$ 過小         & $\eta\!\uparrow$（$0.1\!\to\!0.2$）\\
      ゆっくり発散       & ワインドアップ      & $T_i\!\uparrow$ \\
      \hline
    \end{tabular}
  \end{table}

  \vspace{0.5em}

  \begin{block}{次のステップ}
    Lesson 7 でステップ応答を計測し、この表を使って調整する
  \end{block}
\end{frame}

% ------------------------------------------------------------------
% Slide 7: Recommended Gains (Ti/Td form)
% ------------------------------------------------------------------
\begin{frame}{推奨ゲイン / Recommended Gains}
  \begin{table}
    \centering
    \begin{tabular}{lcccc}
      \hline
      \textbf{軸} & \textbf{$K_p$} & \textbf{$T_i$ [s]} & \textbf{$T_d$ [s]} & \textbf{$\eta$} \\
      \hline
      Roll  & 0.5 & 1.67 & 0.01  & 0.125 \\
      Pitch & 0.5 & 1.67 & 0.01  & 0.125 \\
      Yaw   & 2.0 & 4.0  & 0.005 & 0.125 \\
      \hline
    \end{tabular}
  \end{table}

  \vspace{0.5em}

  \begin{block}{チューニング手順}
    \begin{enumerate}
      \item $T_i{=}0,\; T_d{=}0$ にして $K_p$ を調整（振動しない最大値）
      \item $T_i$ を大きい値 (5.0) から徐々に下げる
      \item $T_d$ を小さい値 (0.001) から徐々に上げる
    \end{enumerate}
  \end{block}
\end{frame}

% ------------------------------------------------------------------
% Slide 8: Hands-on Step 1 -- Ideal PID (extension of L5)
% ------------------------------------------------------------------
\begin{frame}[fragile]{実習 Step\,1: 理想 PID / Ideal PID}
  \begin{lstlisting}[style=sfcpp, title={student.cpp (roll axis excerpt)}]
float Kp=0.5f, Ti=1.67f, Td=0.01f;
float integral=0, prev_err=0;
// ... inside loop_400Hz(dt) ...
float target = ws::rc_roll() * 1.0f;
float error  = target - ws::gyro_x();
float P = Kp * error;
// I term (trapezoidal)
if (Ti > 0)
    integral += (dt/(2*Ti)) * (error + prev_err);
float I = Kp * integral;
// D term (ideal derivative)
float D = Kp * Td * (error - prev_err) / dt;
prev_err = error;
float roll_out = P + I + D;
  \end{lstlisting}

  \begin{exampleblock}{やってみよう}
    飛ばしてみて $\to$ 高周波振動に注目
  \end{exampleblock}
\end{frame}

% ------------------------------------------------------------------
% Slide 9: Hands-on Step 2 -- Incomplete Derivative Filter
% ------------------------------------------------------------------
\begin{frame}[fragile]{実習 Step\,2: 不完全微分 / Incomplete Derivative}
  {\scriptsize
  \begin{lstlisting}[style=sfcpp, title={student.cpp (D term replacement)}]
// Replace ideal D with incomplete derivative filter
float eta = 0.125f;
float d_filt = 0;  // persistent state
// ... inside loop_400Hz(dt) ...
float alpha = 2*eta*Td / dt;
float a = (alpha - 1) / (alpha + 1);
float b = 2*Td / ((alpha + 1) * dt);
d_filt = a * d_filt + b * (error - prev_err);
float D = Kp * d_filt;  // Filtered!
  \end{lstlisting}
  }

  \begin{exampleblock}{やってみよう}
    再び飛ばして $\to$ 高周波振動が減ったことを確認
  \end{exampleblock}
\end{frame}

% ------------------------------------------------------------------
% Slide 10: Checkpoint (updated)
% ------------------------------------------------------------------
\begin{frame}{チェックポイント / Checkpoint}
  \begin{alertblock}{確認事項}
    \begin{itemize}
      \item[$\square$] 理想微分 vs 不完全微分で高周波振動の違いを確認
      \item[$\square$] 定常偏差がなくなった（P のみと比較）
      \item[$\square$] ゲイン調整表を使って応答を改善できた
    \end{itemize}
  \end{alertblock}

  \vspace{1em}

  \begin{block}{次のレッスン}
    Lesson 7: テレメトリ + ステップ応答 --- データで性能を評価
  \end{block}
\end{frame}

\end{document}
